%
% 13-wurzel.tex -- Die Wurzelfunktion
%
% (c) 2026 Prof Dr Andreas Müller
%
\subsection{Die Wurzelfunktion}
Die Wurzelfunktion ist eine Lösung der algebraischen Gleichung $y^2=x$.
Durch Ableiten nach $x$ wird daraus die nichtlineare Differentialgleichung
\begin{equation}
2yy'=1.
\label{buch:differentialgleichungen:potenzreihen:eqn:wurzeldgl}
\end{equation}
Mit der Funktion
\[
F(x,y,y')
=
2yy'-1
\]
kann die Differentialgleichung
\eqref{buch:differentialgleichungen:potenzreihen:eqn:wurzeldgl}
in der Form
\eqref{buch:differentialgleichungen:potenzreihenansatz:eqn:dgl}
geschrieben werden.
Die Funktion $F$ ist eine Polynom in den Variablen $y$ und $y'$ und
ist damit insbesondere auch für komplexe Argumente definiert.
\eqref{buch:differentialgleichungen:potenzreihen:eqn:wurzeldgl}
ist eine komplexe Differentialgleichung für die Wurzelfunktion
und sollte daher eine komplexe Potenzreihe für die Wurzelfunktion
zu liefern in der Lage sein.

%
% Potenzreihenlösung für die Wurzelfunktion
%
\subsubsection{Potenzreihenlösung für die Wurzelfunktion}
Wir schreiben die Lösung als
\[
y(x)
=
\sum_{k=0}^\infty a_kx^k
\qquad\Rightarrow\qquad
y'(x)
=
\sum_{k=1}^\infty ka_kx^{k-1}
=
\sum_{l=0}^\infty (l+1)a_{l+1}x^l
\]
und setzen in die Differentialgleichung
\eqref{buch:differentialgleichungen:potenzreihen:eqn:wurzeldgl}
ein.
Es entsteht die Gleichung
\begin{align}
\frac12
&=
\sum_{k=0}^\infty \sum_{l=0}^\infty a_k(l+1)a_{l+1}x^{k+l}
=
\sum_{n=0}^\infty
\biggl(
\sum_{l=0}^n (l+1)a_{n-l}a_{l+1}
\biggr)
x^n.
\end{align}
Durch Koeffizientenvergleich können wir jetzt die Koeffizienten $a_k$
bestimmen.


Der Koeffizient für $n=0$ ist $a_0a_1=\frac12$ oder $a_1=\frac{1}{2a_0}$.
Die Koeffizienten von $x^n$ mit $n>0$ müssen alle verschwinden, es muss
also
\begin{align*}
0
&=
\sum_{l=0}^n (l+1)a_{n-l}a_{l+1}
&&\Rightarrow&
(n+1)a_0a_{n+1}
&=
-
\sum_{l=0}^{n-1} (l+1)a_{n-l}a_{l+1}
\\
&
&&\Rightarrow&
a_{n+1}
&=
\frac{1}{(n+1)a_0}
\sum_{l=0}^{n-1} (l+1)a_{n-l}a_{l+1}.
\end{align*}
Dies ist eine Rekursionsformel, die den Koeffizienten $a_{n+1}$ aus
den Koeffizienten $a_0,\dots,a_n$ bestimmt.
Alle Koeffizienten hängen also vom Wert $a_0$ ab.

%
% Der Spezialfall a_0=1
%
\subsubsection{Der Spezialfall $a_0=1$}
Um die Rechnung zu vereinfachen, berechnen wir die Koeffizienten nur für 
den Fall $a_0=1$, also für den Fall einer Wurzelfunktion, die für $x=0$
den Wert $1$ ergibt.
Um die Koeffizienten besser auseinanderzuhalten, bezeichnen wir die
Koeffizienten für diesen Spezialfall mit $b_k$.
%
% fig-wurzel.tex
%
% (c) 2026 Prof Dr Andreas Müller
%
\begin{figure}
\centering
\includegraphics{chapters/050-differentialgleichungen/images/wurzel.pdf}
\caption{Potenzreihenapproximation der Wurzelfunktionh $y(x)=\!\sqrt{1+x}$
durch ein Polynom, welches aus einem Potenzreihenansatz für die Lösung
der Differentialgleichung $yy'=\frac12$ mit der Anfangsbedingung
$y(0)=1$ gefunden wurde.
\label{buch:differentialgleichungen:fig:wurzel}}
\end{figure}
%
Wir erwarten also eine Potenzreihenentwicklung für die Funktion
$y(x)=\!\sqrt{1+x}$.
Die Rechnung ergibt
\begin{align*}
n&=0: &&&
1 &=
b_0{\color{darkred}b_1}
\\
&&&\Rightarrow&
b_1
&=
\frac{1}{2b_0}
&&=
\frac12
\\
n&=1: &&&
0 &=
b_1^2+2b_0{\color{darkred}b_2}
\\
&&&\Rightarrow&
{\color{darkred}b_2}
&= 
-\frac{b_1^2}{2b_0}
&&=
-\frac{1}{8}
\\
n&=2: &&&
0 &=
b_2b_1+2b_1b_2+3b_0{\color{darkred}b_3}
\\
&&&\Rightarrow&
{\color{darkred}b_3}
&=
-\frac{1}{3b_0}(3b_1b_2)
&&=
\frac{1}{16}
\\
n&=3: &&&
0 &=
b_3b_1+2b_2^2+3b_1b_3+4b_0{\color{darkred}b_4}
\\
&&&\Rightarrow&
{\color{darkred}b_4}
&=
-
\frac{1}{4b_0}(4b_3b_1+2b_2^2)
&&=
-\frac{5}{128}
\\
n&=4: &&&
0 &=
b_4b_1+2b_3b_2+3b_2b_3+4b_1b_4+5b_0{\color{darkred}b_5}
\\
&&&\Rightarrow&
{\color{darkred}b_5}
&=
-\frac{1}{5b_0}(5b_1b_4+5b_2b_3)
&&=
\frac{7}{256}
\end{align*}
Die Lösung hat also die Potenzreihe
\begin{align}
\!\sqrt{1+x}
&=
\frac12
-\frac{1}{8}x
+\frac{1}{16}x^2
-\frac{5}{128}x^3
+\frac{7}{256}x^4
-\dots
\label{buch:differentialgleichungen:potenzreihen:wurzelapprox}
\end{align}
Die Abbildung~\ref{buch:differentialgleichungen:fig:wurzel}
zeigt die exakte Lösung $\!\sqrt{1+x}$ im Vergleich zur Approximation
\eqref{buch:differentialgleichungen:potenzreihen:wurzelapprox}.

%
% Der allgemeine Fall
%
\subsubsection{Der allgemeine Fall}
Die Differentialgleichung
\eqref{buch:differentialgleichungen:potenzreihen:eqn:wurzeldgl}
hat als allgemeine Lösung die Funktion $y(x)=\!\sqrt{x-C}$.
Die Wahl $a_0=1$ entspricht daher dem Wert $C=-1$.
Für beliebiges $a_0$ findet man
\[
y(x)
=
\!\sqrt{x+a_0^2}
=
\frac{1}{a_0}
\sqrt{1+\frac{x\mathstrut}{a_0^2\mathstrut}}
=
\frac{1}{a_0}
\sum_{k=0}^\infty b_k\biggl(\frac{x\mathstrut}{{a_0}^2\mathstrut}\biggr)^k
\]
als Lösung.
Daraus können wir ablesen, dass
\[
a_k
=
\frac{b_k}{a_0^{2k+1}}
\]
ist.

%
% Die binomische Reihe
%
\subsubsection{Die binomische Reihe}
Für die Funktion $f(x) = (1+x)^\alpha$ ist auch die binomische Formel
\begin{align*}
(1+x)^\alpha
&=
\sum_{k=0}^\infty \binom{\alpha}{k}x^k
=
\sum_{k=0}^\infty \frac{\alpha(\alpha-1)\cdot\ldots\cdot (\alpha-k+1)}{k!}x^k
\end{align*}
bekannt
Im Spezialfall $\alpha=\frac12$ kann man die Binomialkoeffizienten
explizit ausrechnen.
Dazu erweitern wir den Bruch
\begin{align}
\binom{\frac12}{k}
&=
\frac{\frac12\cdot(-\frac12)\cdot(-\frac32)\cdot\ldots\cdot(-\frac{2k-3}{2})}{k!}
=
(-1)^{k-1}
\frac{1\cdot 3 \cdot\ldots\cdot(2k-3)}{2^k\cdot k!}
\intertext{mit $2^k\cdot k!=2\cdot 4 \cdot \ldots\cdot (2k)$ und erhalten}
&=
(-1)^{k-1}
\frac{1\cdot 2\cdot 3 \cdot 4 \cdot\ldots\cdot(2k-3)\cdot(2k-2)\cdot 2k}{(2^k\cdot k!)^2}
\notag
\\
&=
(-1)^{k-1}
\frac{(2k)!}{2^{2k}\cdot (k!)^2\cdot(2k-1)}
=
(-1)^{k-1}
\frac{1}{2^{2k}}
\binom{2k}{k}.
\label{buch:differentialgleichungen:potenzreihen:eqn:koeffizient}
\end{align}
Tatsächlich bekommt
\eqref{buch:differentialgleichungen:potenzreihen:eqn:koeffizient}
für die ersten paar $k$ die Werte der Koeffizienten der
Tabelle~\ref{buch:differentialgleichungen:tab:wurzel}.
%
% tab-wurzel.tex
%
% (c) 2026 Prof Dr Andreas Müller
%
\begin{table}
\centering
\renewcommand{\arraystretch}{1.4}
\begin{tabular}{|>{$}c<{$}|>{$}c<{$}>{$}c<{$}>{$}c<{$}>{$}c<{$}>{$}c<{$}>{$}c<{$}>{$}c<{$}>{$}c<{$}>{$}c<{$}>{$}c<{$}>{$}c<{$}|}
\hline
k & 0 & 1 & 2 & 3 & 4 & 5 & 6 & 7 & 8 & 9 & 10\\ 
\hline
b_k & 1 & \frac12 & -\frac18 & \frac1{16} & -\frac{5}{128} & \frac{7}{256} &
-\frac{21}{1024} & \frac{33}{2048} & -\frac{429}{32768} & -\frac{715}{65536} &
-\frac{2431}{262144}
\\[2pt]
\hline
\end{tabular}
\caption{Koeffizienten der Binomialreihe für $\!\sqrt{1+x}=(1+x)^{\frac12}$
\label{buch:differentialgleichungen:tab:wurzel}}
\end{table}
%
Sie stimmen mit den Koeffizienten in
\eqref{buch:differentialgleichungen:potenzreihen:wurzelapprox}
überein.

